%% fig_decision_tree.tex  --  decision tree for choosing a tracker.
%% %% fig_decision_tree.tex  --  decision tree for choosing a tracker.
%% %% fig_decision_tree.tex  --  decision tree for choosing a tracker.
%% %% fig_decision_tree.tex  --  decision tree for choosing a tracker.
%% \input{figs/fig_decision_tree.tex} inside a figure environment; needs only
%% tikz + shapes.geometric + positioning (all in the main.tex preamble). Root at the top, questions
%% in diamonds, methods in shaded leaves, one line of evidence (mean joint
%% NMI, DynMux vs multislice) under each leaf. Numbers: v1.3.1 rerun, 2026-09-21
%% (multislice n_starts = 10).
\begin{tikzpicture}[
  font=\small, >=stealth,
  q/.style={diamond, aspect=2.2, draw=black!70, thick, fill=white,
            text width=40mm, align=center, inner sep=0pt, font=\small},
  leaf/.style={rectangle, rounded corners=3pt, draw=black!70, thick, fill=black!7,
               text width=34mm, align=center, minimum height=10mm, inner sep=3pt,
               font=\small\bfseries},
  ev/.style={text width=40mm, align=center, font=\scriptsize\color{black!65}, inner sep=1pt},
  arr/.style={->, thick, draw=black!70},
  lab/.style={font=\scriptsize\itshape, fill=white, inner sep=1.5pt}
]

%% level 0
\node[q] (q1) at (0, 0) {Keep each layer's partition exactly as detected, labels aligned afterwards?};

%% level 1
\node[leaf] (hung) at (-4.6, -3.0) {Hungarian label matching};
\node[ev, below=1mm of hung] {procedural guarantee only; never the most accurate (Table~2: 0.28 / 0.32 / 0.36 / 0.10)};
\node[q] (q2) at (3.2, -3.0) {Does the node set or the community set change over time?};

%% level 2
\node[leaf] (dyn) at (-0.6, -6.2) {DynMux\\{\normalfont\small(Jaccard coupling)}};
\node[ev, below=1mm of dyn] {births \& deaths 0.37 vs 0.31\\node turnover 0.36 vs 0.28\\recurrence, period unknown 0.28 vs 0.10};
\node[q] (q3) at (6.4, -6.2) {Nodes persist and mostly keep their community; per-layer signal weak?};

%% level 3
\node[leaf] (ms) at (3.0, -9.4) {Multislice modularity\\{\normalfont\small(identity ties, $\omega = 1$)}};
\node[ev, below=1mm of ms] {gradual switching 0.53 vs 0.28\\core/periphery 0.86 vs 0.46\\long series 0.69 vs 0.23; all splits/merges};
\node[leaf, fill=black!14] (both) at (9.8, -9.4) {Fit both; keep the higher stability score $s$};
\node[ev, below=1mm of both] {picks the more accurate method in 82\% of series (0.47 vs 0.35 / 0.44; oracle 0.48)};

%% edges
\draw[arr] (q1) -- node[lab, above left] {yes} (hung);
\draw[arr] (q1) -- node[lab, above right] {no} (q2);
\draw[arr] (q2) -- node[lab, above left] {yes} (dyn);
\draw[arr] (q2) -- node[lab, above right] {no} (q3);
\draw[arr] (q3) -- node[lab, above left] {yes} (ms);
\draw[arr] (q3) -- node[lab, above right] {unsure} (both);

%% examples under each question (what the question means in practice)
\node[ev, text width=30mm, right=2mm of q2.east, align=left, font=\scriptsize\itshape\color{black!65}] {states enter and exit; blocs form and dissolve; recurrence on an unknown period};
\node[ev, text width=30mm, right=2mm of q3.east, align=left, font=\scriptsize\itshape\color{black!65}] {gradual switching; stable core, rotating periphery; long series; splits and merges};
\end{tikzpicture}
 inside a figure environment; needs only
%% tikz + shapes.geometric + positioning (all in the main.tex preamble). Root at the top, questions
%% in diamonds, methods in shaded leaves, one line of evidence (mean joint
%% NMI, DynMux vs multislice) under each leaf. Numbers: v1.3.1 rerun, 2026-09-21
%% (multislice n_starts = 10).
\begin{tikzpicture}[
  font=\small, >=stealth,
  q/.style={diamond, aspect=2.2, draw=black!70, thick, fill=white,
            text width=40mm, align=center, inner sep=0pt, font=\small},
  leaf/.style={rectangle, rounded corners=3pt, draw=black!70, thick, fill=black!7,
               text width=34mm, align=center, minimum height=10mm, inner sep=3pt,
               font=\small\bfseries},
  ev/.style={text width=40mm, align=center, font=\scriptsize\color{black!65}, inner sep=1pt},
  arr/.style={->, thick, draw=black!70},
  lab/.style={font=\scriptsize\itshape, fill=white, inner sep=1.5pt}
]

%% level 0
\node[q] (q1) at (0, 0) {Keep each layer's partition exactly as detected, labels aligned afterwards?};

%% level 1
\node[leaf] (hung) at (-4.6, -3.0) {Hungarian label matching};
\node[ev, below=1mm of hung] {procedural guarantee only; never the most accurate (Table~2: 0.28 / 0.32 / 0.36 / 0.10)};
\node[q] (q2) at (3.2, -3.0) {Does the node set or the community set change over time?};

%% level 2
\node[leaf] (dyn) at (-0.6, -6.2) {DynMux\\{\normalfont\small(Jaccard coupling)}};
\node[ev, below=1mm of dyn] {births \& deaths 0.37 vs 0.31\\node turnover 0.36 vs 0.28\\recurrence, period unknown 0.28 vs 0.10};
\node[q] (q3) at (6.4, -6.2) {Nodes persist and mostly keep their community; per-layer signal weak?};

%% level 3
\node[leaf] (ms) at (3.0, -9.4) {Multislice modularity\\{\normalfont\small(identity ties, $\omega = 1$)}};
\node[ev, below=1mm of ms] {gradual switching 0.53 vs 0.28\\core/periphery 0.86 vs 0.46\\long series 0.69 vs 0.23; all splits/merges};
\node[leaf, fill=black!14] (both) at (9.8, -9.4) {Fit both; keep the higher stability score $s$};
\node[ev, below=1mm of both] {picks the more accurate method in 82\% of series (0.47 vs 0.35 / 0.44; oracle 0.48)};

%% edges
\draw[arr] (q1) -- node[lab, above left] {yes} (hung);
\draw[arr] (q1) -- node[lab, above right] {no} (q2);
\draw[arr] (q2) -- node[lab, above left] {yes} (dyn);
\draw[arr] (q2) -- node[lab, above right] {no} (q3);
\draw[arr] (q3) -- node[lab, above left] {yes} (ms);
\draw[arr] (q3) -- node[lab, above right] {unsure} (both);

%% examples under each question (what the question means in practice)
\node[ev, text width=30mm, right=2mm of q2.east, align=left, font=\scriptsize\itshape\color{black!65}] {states enter and exit; blocs form and dissolve; recurrence on an unknown period};
\node[ev, text width=30mm, right=2mm of q3.east, align=left, font=\scriptsize\itshape\color{black!65}] {gradual switching; stable core, rotating periphery; long series; splits and merges};
\end{tikzpicture}
 inside a figure environment; needs only
%% tikz + shapes.geometric + positioning (all in the main.tex preamble). Root at the top, questions
%% in diamonds, methods in shaded leaves, one line of evidence (mean joint
%% NMI, DynMux vs multislice) under each leaf. Numbers: v1.3.1 rerun, 2026-09-21
%% (multislice n_starts = 10).
\begin{tikzpicture}[
  font=\small, >=stealth,
  q/.style={diamond, aspect=2.2, draw=black!70, thick, fill=white,
            text width=40mm, align=center, inner sep=0pt, font=\small},
  leaf/.style={rectangle, rounded corners=3pt, draw=black!70, thick, fill=black!7,
               text width=34mm, align=center, minimum height=10mm, inner sep=3pt,
               font=\small\bfseries},
  ev/.style={text width=40mm, align=center, font=\scriptsize\color{black!65}, inner sep=1pt},
  arr/.style={->, thick, draw=black!70},
  lab/.style={font=\scriptsize\itshape, fill=white, inner sep=1.5pt}
]

%% level 0
\node[q] (q1) at (0, 0) {Keep each layer's partition exactly as detected, labels aligned afterwards?};

%% level 1
\node[leaf] (hung) at (-4.6, -3.0) {Hungarian label matching};
\node[ev, below=1mm of hung] {procedural guarantee only; never the most accurate (Table~2: 0.28 / 0.32 / 0.36 / 0.10)};
\node[q] (q2) at (3.2, -3.0) {Does the node set or the community set change over time?};

%% level 2
\node[leaf] (dyn) at (-0.6, -6.2) {DynMux\\{\normalfont\small(Jaccard coupling)}};
\node[ev, below=1mm of dyn] {births \& deaths 0.37 vs 0.31\\node turnover 0.36 vs 0.28\\recurrence, period unknown 0.28 vs 0.10};
\node[q] (q3) at (6.4, -6.2) {Nodes persist and mostly keep their community; per-layer signal weak?};

%% level 3
\node[leaf] (ms) at (3.0, -9.4) {Multislice modularity\\{\normalfont\small(identity ties, $\omega = 1$)}};
\node[ev, below=1mm of ms] {gradual switching 0.53 vs 0.28\\core/periphery 0.86 vs 0.46\\long series 0.69 vs 0.23; all splits/merges};
\node[leaf, fill=black!14] (both) at (9.8, -9.4) {Fit both; keep the higher stability score $s$};
\node[ev, below=1mm of both] {picks the more accurate method in 82\% of series (0.47 vs 0.35 / 0.44; oracle 0.48)};

%% edges
\draw[arr] (q1) -- node[lab, above left] {yes} (hung);
\draw[arr] (q1) -- node[lab, above right] {no} (q2);
\draw[arr] (q2) -- node[lab, above left] {yes} (dyn);
\draw[arr] (q2) -- node[lab, above right] {no} (q3);
\draw[arr] (q3) -- node[lab, above left] {yes} (ms);
\draw[arr] (q3) -- node[lab, above right] {unsure} (both);

%% examples under each question (what the question means in practice)
\node[ev, text width=30mm, right=2mm of q2.east, align=left, font=\scriptsize\itshape\color{black!65}] {states enter and exit; blocs form and dissolve; recurrence on an unknown period};
\node[ev, text width=30mm, right=2mm of q3.east, align=left, font=\scriptsize\itshape\color{black!65}] {gradual switching; stable core, rotating periphery; long series; splits and merges};
\end{tikzpicture}
 inside a figure environment; needs only
%% tikz + shapes.geometric + positioning (all in the main.tex preamble). Root at the top, questions
%% in diamonds, methods in shaded leaves, one line of evidence (mean joint
%% NMI, DynMux vs multislice) under each leaf. Numbers: v1.3.1 rerun, 2026-09-21
%% (multislice n_starts = 10).
\begin{tikzpicture}[
  font=\small, >=stealth,
  q/.style={diamond, aspect=2.2, draw=black!70, thick, fill=white,
            text width=40mm, align=center, inner sep=0pt, font=\small},
  leaf/.style={rectangle, rounded corners=3pt, draw=black!70, thick, fill=black!7,
               text width=34mm, align=center, minimum height=10mm, inner sep=3pt,
               font=\small\bfseries},
  ev/.style={text width=40mm, align=center, font=\scriptsize\color{black!65}, inner sep=1pt},
  arr/.style={->, thick, draw=black!70},
  lab/.style={font=\scriptsize\itshape, fill=white, inner sep=1.5pt}
]

%% level 0
\node[q] (q1) at (0, 0) {Keep each layer's partition exactly as detected, labels aligned afterwards?};

%% level 1
\node[leaf] (hung) at (-4.6, -3.0) {Hungarian label matching};
\node[ev, below=1mm of hung] {procedural guarantee only; never the most accurate (Table~2: 0.28 / 0.32 / 0.36 / 0.10)};
\node[q] (q2) at (3.2, -3.0) {Does the node set or the community set change over time?};

%% level 2
\node[leaf] (dyn) at (-0.6, -6.2) {DynMux\\{\normalfont\small(Jaccard coupling)}};
\node[ev, below=1mm of dyn] {births \& deaths 0.37 vs 0.31\\node turnover 0.36 vs 0.28\\recurrence, period unknown 0.28 vs 0.10};
\node[q] (q3) at (6.4, -6.2) {Nodes persist and mostly keep their community; per-layer signal weak?};

%% level 3
\node[leaf] (ms) at (3.0, -9.4) {Multislice modularity\\{\normalfont\small(identity ties, $\omega = 1$)}};
\node[ev, below=1mm of ms] {gradual switching 0.53 vs 0.28\\core/periphery 0.86 vs 0.46\\long series 0.69 vs 0.23; all splits/merges};
\node[leaf, fill=black!14] (both) at (9.8, -9.4) {Fit both; keep the higher stability score $s$};
\node[ev, below=1mm of both] {picks the more accurate method in 82\% of series (0.47 vs 0.35 / 0.44; oracle 0.48)};

%% edges
\draw[arr] (q1) -- node[lab, above left] {yes} (hung);
\draw[arr] (q1) -- node[lab, above right] {no} (q2);
\draw[arr] (q2) -- node[lab, above left] {yes} (dyn);
\draw[arr] (q2) -- node[lab, above right] {no} (q3);
\draw[arr] (q3) -- node[lab, above left] {yes} (ms);
\draw[arr] (q3) -- node[lab, above right] {unsure} (both);

%% examples under each question (what the question means in practice)
\node[ev, text width=30mm, right=2mm of q2.east, align=left, font=\scriptsize\itshape\color{black!65}] {states enter and exit; blocs form and dissolve; recurrence on an unknown period};
\node[ev, text width=30mm, right=2mm of q3.east, align=left, font=\scriptsize\itshape\color{black!65}] {gradual switching; stable core, rotating periphery; long series; splits and merges};
\end{tikzpicture}
